\section{Grassmann--Cayley meet and projective duality}
\label{sec:molecular-meet}

This chapter develops the dual half of the Grassmann--Cayley algebra of
Katoh--Tanigawa~\cite[\S2.1]{katohTanigawa2011}: the \emph{meet} (regressive
product) and the projective-duality dictionary it rests on, built on the
homogeneous-coordinate carrier of \cref{sec:molecular-homog}. That chapter
built the \emph{join} (the progressive product $A \vee B$, the symbolic
exterior product on $\bigwedge[\R] \R^{k+2}$) and a coordinatized Plücker
bridge; here we build its projective dual. The meet is the device the
panel-hinge realizations of \cref{sec:molecular-algebraic-induction}, the
cycle realization (\cref{lem:cycle-realization}, KT Lemma~5.4, due to
Crapo--Whiteley~\cite{crapoWhiteley1982}), and Katoh--Tanigawa's later
projective-invariance argument all build on.

The construction is \textbf{metric-free}: projective geometry needs no
inner product, only the top-power volume form (an orientation). Write
$N = k+2$ and $V = \R^{N}$ for the homogeneous-coordinate carrier; the top
graded piece $\bigwedge^{N} V$ is free of rank $\binom{N}{N} = 1$, hence
canonically $\cong \R$. In order: the volume-form isomorphism
(\cref{def:meet-top-equiv}) gives the perfect wedge pairing below its
codomain; the duality-pairing isomorphism (\cref{def:meet-pairing-dual}) is
the projective-duality dictionary entry; the complement isomorphism
(\cref{def:meet-complement-iso}), built from that pairing, is the
chapter's central construction; and the meet itself (\cref{def:meet}) is
the join conjugated by the complement isomorphism. The chapter closes with
the point-join $\leftrightarrow$ panel-meet duality
(\cref{lem:case-III-claim612-line-in-panel-union}, KT eq.~(6.45)), the
identity the $d = 3$ Case-III realization of
\cref{sec:molecular-algebraic-induction} is built on.

\begin{definition}[Top-power volume form $\bigwedge^{N} V \cong \R$]
  \label{def:meet-top-equiv}
  \lean{CombinatorialRigidity.Molecular.screwAlgebraTopEquiv,
    exteriorPower.topEquiv}
  \leanok
  The $N$-th exterior power of a free module of rank $N$ is free of rank
  $\binom{N}{N} = 1$, hence canonically isomorphic to the base ring. On
  the standard carrier $V = \R^{N}$ this is the \emph{volume
  form}: the linear isomorphism
  \[
    \mathrm{topEquiv} : \bigwedge^{N} V \xrightarrow{\ \sim\ } \R,
  \]
  sending the wedge $\mathbf{e}_1 \wedge \cdots \wedge \mathbf{e}_N$ of all
  standard basis vectors to $1$. It carries an orientation, not a metric:
  the meet is built from it without any inner-product structure.
\end{definition}
\begin{proof}
  \leanok
  Built from the standard-basis top-power basis
  $\mathrm{Module.Basis.exteriorPower}\,(\mathrm{Pi.basisFun}\ \R\ (\mathrm{Fin}\ N))$,
  whose index set $\mathrm{powersetCard}\,(\mathrm{Fin}\ N)\ N$ is a
  singleton (the full set is the unique $N$-element subset of an
  $N$-element type), composed with the unique-index evaluation
  $\mathrm{LinearEquiv.funUnique}$. The general fact (over any commutative
  ring, on $\mathrm{Fin}\ n \to R$) is the mirror lemma
  $\mathrm{exteriorPower.topEquiv}$; $\mathrm{screwAlgebraTopEquiv}$ is the
  $N = k+2$ specialization on the screw-algebra carrier of
  \cref{sec:molecular-homog}.
\end{proof}

\begin{definition}[Exterior-power duality iso $\bigwedge^{j} V^{*} \cong (\bigwedge^{j} V)^{*}$]
  \label{def:meet-pairing-dual}
  \lean{CombinatorialRigidity.Molecular.screwAlgebraPairingDualEquiv,
    exteriorPower.pairingDualEquiv}
  \leanok
  \uses{def:meet-top-equiv}
  The natural pairing map $\bigwedge^{j} V^{*} \to (\bigwedge^{j} V)^{*}$,
  sending a dual-exterior-power element to the functional it induces on
  $\bigwedge^{j} V$, is for free finite $V$ an isomorphism
  \[
    \mathrm{pairingDualEquiv} : \bigwedge^{j} V^{*}
      \xrightarrow{\ \sim\ } (\bigwedge^{j} V)^{*}.
  \]
\end{definition}
\begin{proof}
  \leanok
  Fix an ordered basis $b$ of $V$. The map $\mathrm{pairingDual}$ carries
  the exterior-power basis $\bigwedge^{j} b^{*}$ built from the dual basis
  onto the dual basis of the exterior-power basis $\bigwedge^{j} b$: the
  identification $b^{*} = b.\mathrm{coord}$ together with the coordinate
  computation $\mathrm{exteriorPower.basis\_coord}$ shows the image of each
  basis vector is the corresponding coordinate form. Hence
  $\mathrm{pairingDual}$ sends a basis to a basis and is an isomorphism;
  the mirror lemma $\mathrm{exteriorPower.pairingDualEquiv}$ packages this
  as a $\mathrm{Basis.equiv}$ and $\mathrm{coe\_pairingDualEquiv}$ identifies
  it with $\mathrm{pairingDual}$ in place.
  $\mathrm{screwAlgebraPairingDualEquiv}$ is the specialization at the
  standard basis of the screw-algebra carrier.
\end{proof}

\begin{definition}[Complement iso $\bigwedge^{j} V \cong \bigwedge^{N-j} V$]
  \label{def:meet-complement-iso}
  \lean{CombinatorialRigidity.Molecular.complementIso}
  \leanok
  \uses{def:meet-top-equiv}
  The perfect wedge pairing
  $\bigwedge^{j} V \times \bigwedge^{N-j} V \to \bigwedge^{N} V
  \cong \R$ (the join of \cref{sec:molecular-homog} followed by the
  volume form of \cref{def:meet-top-equiv}) is nondegenerate for free $V$,
  inducing a linear isomorphism
  $\mathrm{complementIso} : \bigwedge^{j} V \cong \bigwedge^{N-j} V$.
\end{definition}
\begin{proof}
  \leanok
  Nondegeneracy is checked on the standard exterior-power basis, whose
  vectors are indexed by the $j$- and $(N-j)$-element subsets $S$ and $T$ of
  $\{1, \dots, N\}$. The pairing of $\mathbf{e}_S$ and $\mathbf{e}_T$ is the
  volume form applied to the join $\mathbf{e}_S \vee \mathbf{e}_T$, i.e.\ the
  extensor of the concatenated index family. If $S$ and $T$ overlap (which,
  for complementary cardinalities, is exactly $T \ne S^{c}$) the family
  repeats a basis vector and the extensor vanishes by alternation. If
  $T = S^{c}$ the family is a permutation of all $N$ standard basis vectors,
  so it is linearly independent and its extensor is nonzero. Hence the pairing
  matrix on the standard basis is a signed-permutation matrix, so the pairing
  $\mathrm{wedgePairing} : \bigwedge^{j} V \to (\bigwedge^{N-j} V)^{*}$ is
  injective; as $\bigwedge^{j} V$ and $\bigwedge^{N-j} V$ have equal dimension
  $\binom{N}{j} = \binom{N}{N-j}$, injectivity makes it a linear isomorphism
  onto the dual. Post-composing with the dual-evaluation iso of
  $\bigwedge^{N-j} V$ lands $\mathrm{complementIso}$ in $\bigwedge^{N-j} V$
  itself. The exact grade-swap sign is not needed for the isomorphism and is
  deferred to where an oriented meet consumes it.
\end{proof}

\begin{lemma}[Defining wedge-pairing property of the complement iso]
  \label{lem:complement-iso-toDual}
  \lean{CombinatorialRigidity.Molecular.complementIso_toDual}
  \leanok
  \uses{def:meet-complement-iso, def:meet-top-equiv}
  Pairing $\mathrm{complementIso}\,X \in \bigwedge^{N-j} V$ against any
  $B \in \bigwedge^{N-j} V$ through the standard exterior-power basis's
  $\mathrm{toDual}$ reproduces the wedge pairing of $X$ with $B$:
  \[
    \langle \mathrm{complementIso}\,X,\ B \rangle
      = \mathrm{wedgePairing}(X, B)
      = \mathrm{vol}\,(X \vee B).
  \]
  This is the characterizing identity of the complement iso.
\end{lemma}
\begin{proof}
  \leanok
  By construction $\mathrm{complementIso} = (\mathrm{wedgePairing}\text{ as an
  equiv}) \mathbin{\circ} \mathrm{toDualEquiv}^{-1}$; applying $\mathrm{toDualEquiv}$
  undoes the second factor (\texttt{LinearEquiv.apply\_symm\_apply}) and leaves
  $\mathrm{wedgePairing}\,X$.
\end{proof}

\begin{lemma}[Complement of a decomposable lies in the wedge-orthogonal complement]
  \label{lem:complement-iso-decomposable-wedge-perp}
  \lean{CombinatorialRigidity.Molecular.complementIso_toDual_eq_zero_of_wedgeProd_eq_zero,
        CombinatorialRigidity.Molecular.wedgeProd_extensor_eq_zero_of_shared_vector,
        CombinatorialRigidity.Molecular.complementIso_toDual_extensor_eq_zero_of_shared_vector}
  \leanok
  \uses{def:meet-complement-iso, lem:complement-iso-toDual, def:join, def:extensor}
  The first of two steps toward the point-join $\leftrightarrow$ panel-meet
  duality (\cref{lem:case-III-claim612-line-in-panel-union}). Let $n_u, n'$
  be two independent vectors of $\R^4$ --- in the application below, two
  normals of a panel $\Pi(u)$ in the sense of
  \cref{sec:molecular-algebraic-induction} --- and $W = \{n_u, n'\}^{\perp}$.
  The complement $\mathrm{complementIso}(n_u \wedge n')$ of the grade-$2$
  decomposable $n_u \wedge n'$ lies in $\bigwedge^{2} W$, in the operational
  dual sense that it is annihilated by every $2$-extensor sharing a vector
  with $n_u \wedge n'$: through the dictionary entry
  \cref{lem:complement-iso-toDual}, a vanishing wedge has vanishing
  complement pairing ($X \vee B = 0 \Rightarrow \langle
  \mathrm{complementIso}\,X, B\rangle = \mathrm{vol}(X \vee B) = 0$), and at
  $d = 3$ ($N = 4$) the wedge of two $2$-extensors sharing a vector vanishes
  by alternation (the concatenated family repeats the shared vector).
\end{lemma}
\begin{proof}
  \leanok
  The dictionary half is immediate from \cref{lem:complement-iso-toDual} and linearity of the volume
  form ($\mathrm{vol}\,0 = 0$). The concrete half writes the underlying element of the graded wedge
  as the join $\mathrm{extensor}\,n \vee \mathrm{extensor}\,c = \mathrm{extensor}(\mathrm{append}\,n\,
  c)$ (\cref{def:join}); a shared vector makes the appended family non-injective, so the extensor
  vanishes by the alternating law (\cref{def:extensor}).
\end{proof}

\begin{lemma}[The exterior square of a $2$-dimensional space is a line]
  \label{lem:complement-iso-line-one-dim}
  \lean{CombinatorialRigidity.Molecular.finrank_exteriorPower_two_eq_one,
        CombinatorialRigidity.Molecular.exteriorPower_finrank_eq_one_proportional}
  \leanok
  \uses{def:join, def:extensor}
  The second step toward the point-join $\leftrightarrow$ panel-meet duality
  (\cref{lem:case-III-claim612-line-in-panel-union}): for a $2$-dimensional
  space $W$, the exterior square $\bigwedge^{2} W$ is $1$-dimensional,
  \[
    \dim \bigwedge^{2} W = \binom{\dim W}{2} = \binom{2}{2} = 1.
  \]
  Consequently any two members of $\bigwedge^{2} W$, one of them nonzero,
  are scalar multiples of one another.
\end{lemma}
\begin{proof}
  \leanok
  The dimension count is $\mathrm{exteriorPower.finrank\_eq}$ ($\dim \bigwedge^{n} W = \binom{\dim
  W}{n}$) at $\dim W = 2$, $n = 2$: $\binom{2}{2} = 1$. A space of dimension $1$ with a fixed nonzero
  vector $x$ has every vector a multiple of $x$ ($\mathrm{finrank\_eq\_one\_iff\_of\_nonzero}'$).
\end{proof}

\begin{lemma}[The exterior square of a subspace includes into $\bigwedge^{2} \R^4$ injectively]
  \label{lem:complement-iso-exterior-map-injective}
  \lean{CombinatorialRigidity.Molecular.exteriorPower_map_subtype_injective}
  \leanok
  \uses{def:meet-complement-iso}
  A further step toward the point-join $\leftrightarrow$ panel-meet duality
  (\cref{lem:case-III-claim612-line-in-panel-union}), on top of
  \cref{lem:complement-iso-decomposable-wedge-perp} and
  \cref{lem:complement-iso-line-one-dim}. For a subspace $W \subseteq \R^4$,
  the exterior-power map $\bigwedge^{2}(W.\mathrm{subtype}) : \bigwedge^{2}
  W \to \bigwedge^{2} \R^4$ induced by the inclusion $W \hookrightarrow \R^4$
  is injective.
\end{lemma}
\begin{proof}
  \leanok
  Over the field $\R$, an injective linear map induces an injective exterior power
  ($\mathrm{exteriorPower.map\_injective\_field}$); the inclusion $W.\mathrm{subtype}$ is injective
  ($\mathrm{Submodule.injective\_subtype}$). (No retraction is needed: the $\mathrm{CommRing}$-level
  injectivity of $\bigwedge^{n} f$ requires an explicit left inverse of $f$, but over a field
  injectivity of $f$ alone suffices.)
\end{proof}

\begin{lemma}[Point-join $\leftrightarrow$ panel-meet duality; KT eq.~(6.45)]
  \label{lem:case-III-claim612-line-in-panel-union}
  \lean{CombinatorialRigidity.Molecular.complementIso_smul_eq_extensor_join,
        CombinatorialRigidity.Molecular.extensor_join_eq_zero_of_complementIso_eq_zero,
        CombinatorialRigidity.Molecular.extensor_join_eq_zero_of_complementIso_eq_zero_dotProduct,
        CombinatorialRigidity.Molecular.panelSupportExtensor_join_eq_zero_of_eq_zero,
        CombinatorialRigidity.Molecular.panelSupportExtensor_add_smul_left_ne_zero_of_join_ne_zero,
        CombinatorialRigidity.Molecular.exists_shear_linearIndependent_pair,
        CombinatorialRigidity.Molecular.PanelHingeFramework.case_III_old_new_blocks_of_line,
        CombinatorialRigidity.Molecular.PanelHingeFramework.case_III_full_family_of_line,
        CombinatorialRigidity.Molecular.PanelHingeFramework.case_III_realization_of_line,
        CombinatorialRigidity.Molecular.exists_line_data_of_homogeneousIncidence,
        CombinatorialRigidity.Molecular.BodyHingeFramework.exists_complementIso_ne_zero_of_homogeneousIncidence,
        CombinatorialRigidity.Molecular.PanelHingeFramework.hasGenericFullRankRealization_of_rigidOn_ofNormals}
  \leanok
  \uses{lem:complement-iso-decomposable-wedge-perp, lem:complement-iso-line-one-dim,
        lem:complement-iso-exterior-map-injective, def:meet-complement-iso, def:join, def:extensor}
  At $d = 3$ ($\bigwedge^{2}\R^4$), let $n_u, n'$ be two independent panel normals of a panel
  $\Pi(u)$ and $p_i, p_j$ two points whose connecting line $L = p_i p_j$ lies in $\Pi(u)$ (each
  point orthogonal to both normals). Then the point-join $\bar{p_i} \vee \bar{p_j} =
  \mathrm{extensor}\,[p_i, p_j]$ is a scalar multiple of the panel-meet $C(L) =
  \mathrm{complementIso}(n_u \wedge n')$ --- both are the Plücker vector of the line $L$, up to the
  projective scale. Consequently a screw functional $r$ annihilating $C(L)$ also annihilates
  $\bar{p_i} \vee \bar{p_j}$. This is the duality Katoh--Tanigawa use implicitly in reading
  eq.~(6.45): the spanning point-joins and the annihilated panel-meets of the lines in a panel
  union $\Pi(a) \cup \Pi(b) \cup \Pi(c)$ are the \emph{same} extensors.
\end{lemma}
\begin{proof}
  \leanok
  Both members lie in the common line $\bigwedge^{2}W$ for $W = \{n_u, n'\}^\perp$ (the
  wedge-orthogonal complement of the two independent normals, of dimension $4 - 2 = 2$): the
  point-join lies there because each of $p_i, p_j$ is orthogonal to both normals, and the
  panel-meet lies there by the decomposable-of-orthogonal-complement fact
  (\cref{def:meet-complement-iso}). Since $\dim W = 2$, $\bigwedge^{2}W$ is
  $1$-dimensional (\cref{lem:complement-iso-line-one-dim}, in the $\bigwedge^{2}\R^4$-internal form
  via the injective pull-back of \cref{lem:complement-iso-exterior-map-injective}), and the
  panel-meet is nonzero (the normals are independent), so it spans the line and the point-join is
  its multiple, $\bar{p_i} \vee \bar{p_j} = \lambda\, C(L)$. The annihilation transfer is then
  immediate: $r(C(L)) = 0$ gives $r(\bar{p_i} \vee \bar{p_j}) = \lambda \cdot r(C(L)) = 0$.
\end{proof}
\begin{fmlnote}
  The declaration list also carries the $d = 3$ Case-III realization built on this duality
  (\cref{sec:molecular-algebraic-induction}): reading the identity in the row-independence test's
  own coordinates (\texttt{panelSupportExtensor\_join\_eq\_zero\_of\_eq\_zero}), transporting it
  along a sheared candidate normal (\texttt{panelSupportExtensor\_add\_smul\_left\_ne\_zero\_of\_join\_ne\_zero},
  \texttt{exists\_shear\_linearIndependent\_pair}), assembling the resulting independent row family
  into a full-rank realization (\texttt{case\_III\_old\_new\_blocks\_of\_line},
  \texttt{case\_III\_full\_family\_of\_line}, \texttt{case\_III\_realization\_of\_line},
  \texttt{exists\_line\_data\_of\_homogeneousIncidence},
  \texttt{exists\_complementIso\_ne\_zero\_of\_homogeneousIncidence}), and upgrading that bare
  realization to a generic one (\texttt{hasGenericFullRankRealization\_of\_rigidOn\_ofNormals}, KT's
  Lemma~5.2 "convert to a nonparallel realization without decreasing rank"). The last of these is a
  general-purpose upgrade, not specific to Case III: it is also what \cref{lem:cycle-realization}
  cites for its own bare-to-generic step.
\end{fmlnote}

\begin{definition}[The regressive product (meet)]
  \label{def:meet}
  \lean{CombinatorialRigidity.Molecular.meet}
  \leanok
  \uses{def:meet-complement-iso}
  The \emph{meet} (regressive product) $A \wedge B$ of two extensors is the
  Grassmann--Cayley dual of the join: writing $*$ for the complement iso
  $\bigwedge^{j} V \to \bigwedge^{N-j} V$ of \cref{def:meet-complement-iso}, the
  meet of $A \in \bigwedge^{N-a} V$ and $B \in \bigwedge^{N-b} V$ (with
  $a + b \le N$) is $*(*A \vee *B)$, landing in $\bigwedge^{N-(a+b)} V$: the
  complements $*A \in \bigwedge^{a} V$ and $*B \in \bigwedge^{b} V$ are joined to a
  $\bigwedge^{a+b} V$ element (the graded product $\mathrm{gradedMul}$, the
  general-grade form of $\mathrm{wedgeProd}$), and a third $*$ returns the result to
  $\bigwedge^{N-(a+b)} V$. Geometrically, when $A$, $B$ are the supporting extensors
  of two subspaces of $V$, the meet is the supporting extensor of their
  codimension-$(a+b)$ intersection; for $a = b = 1$ the meet of two hyperplane
  normals is the supporting $(N-2) = k$-extensor of their codimension-2
  intersection, landing in the screw space $\mathrm{ScrewSpace}\,k$ of
  \cref{sec:molecular-rigidity-matrix}.
\end{definition}
